\documentclass[10pt,a4paper]{article}
\usepackage[utf8]{inputenc}
\usepackage[a4paper,left=3cm,right=2cm]{geometry}
\usepackage{amsmath}
\usepackage[thmmarks, amsmath, thref]{ntheorem}
\usepackage{amsfonts}
\usepackage{amssymb}
\usepackage{fancyhdr}
\pagestyle{fancy}
\usepackage{anysize}
\usepackage{graphicx}
\usepackage{float}
\usepackage{hyperref}
\usepackage{multirow}
\usepackage{enumitem}
\newtheorem{theorem}{Theorem}
\theoremsymbol{\ensuremath{\blacksquare}}
\newtheorem*{solution}{Solución:}
\usepackage{listings}
\usepackage{xcolor}

% Margenes y formalidades

\textwidth 6.25in % Margen Derecho y ancho del texto
\textheight 8.5in
\oddsidemargin 0in % Margen Izquierdo
\headheight 0in % Largo del texto
\leftmargin 1in
\parindent 0em % Salto entre la indentación (Sangría)https://www.overleaf.com/4519127669sdpptkhvsbfn
\topmargin -0.5in % Margen Superior
\parskip 0.5ex % Salto entre parrafos
\baselineskip 1pt

%\lhead{\includegraphics[scale=0.2]{logo.png}} % texto izquierda de la cabecera
%\chead{}                                      % texto centro de la cabecera
%\rhead{\includegraphics[scale=0.2]{FIG/logo-mba-2016.png}} % número de página a la derecha

\renewcommand{\headrulewidth}{0.4pt} % grosor de la línea de la cabecera
\renewcommand{\footrulewidth}{0.4pt} % grosor de la línea del pie

\begin{document}



\pagebreak
\vspace*{4mm} 

\begin{center}
\begin{huge}
\textbf{\\Ejercicios Econometría 3}
\end{Large}\\
%{\large 2023-2}\\
\end{center}

\begin{enumerate}

  \item  Considere los siguientes datos para estimar el modelo
$\displaystyle Y_i=\beta_0+\beta_1X_i+u_i$, donde $X_i$

\begin{center}
\begin{tabular}{|c|c|c|}
\hline
$i$ & $X_i$ & $Y_i$\\
\hline
1 & 1 & 5\\
2 & 3 & 10\\
3 & 5 & 11\\
4 & 7 & 18\\
5 & 9 & 16\\
\hline
\end{tabular}
\end{center}


\begin{enumerate}
  \item Estime $\hat\beta_0$ y $\hat\beta_1$ mediante MCO.
  \item Calcule SST, SSR y SSE
  \item Obtenga el coeficiente $R^2$
\end{enumerate}
\section*{Soluci\'on}

\begin{enumerate}
%-------------------------------------------------------------------
\item \textbf{Estimación de los parámetros}

Los promedios muestrales son
\[
\bar X \;=\;\frac{1+3+5+7+9}{5}=5,
\qquad
\bar Y \;=\;\frac{5+10+11+18+16}{5}=12.
\]

\[
\hat\beta_1
=\frac{\sum_{i=1}^{5}(X_i-\bar X)(Y_i-\bar Y)}
        {\sum_{i=1}^{5}(X_i-\bar X)^2}
  =\frac{60}{40}
  =\frac{3}{2}=1.5,
\qquad
\hat\beta_0
=\bar Y-\hat\beta_1\bar X
 =12-\frac{3}{2}\cdot5
 =\frac{9}{2}=4.5 .
\]

Así, la recta ajustada es
\[
\widehat{Y}_i \;=\; \hat\beta_0+\hat\beta_1X_i=4.5+1.5X_i.
\]

%-------------------------------------------------------------------
\item \textbf{Cálculo de } $SST$, $SSE$ \textbf{y} $SSR$

\[
\begin{array}{|c|c|c|c|c|}
\hline
i & X_i & Y_i & \widehat Y_i = 4.5+1.5X_i & u_i = Y_i-\widehat Y_i \\
\hline
1 & 1 & 5  & 6  & -1 \\
2 & 3 & 10 & 9  & \; 1 \\
3 & 5 & 11 & 12 & -1 \\
4 & 7 & 18 & 15 & \; 3 \\
5 & 9 & 16 & 18 & -2 \\
\hline
\end{array}
\]

\[
\begin{aligned}
SSR &= \sum_{i=1}^{5} u_i^{\,2}
     = (-1)^2 + 1^2 + (-1)^2 + 3^2 + (-2)^2
     = 16,\\[6pt]
SST &= \sum_{i=1}^{5} (Y_i-\bar Y)^2
     = (5-12)^2 + (10-12)^2 + (11-12)^2 + (18-12)^2 + (16-12)^2
     = 106,\\[6pt]
SSE &= SST - SSR = 106 - 16 = 90.
\end{aligned}
\]

%-------------------------------------------------------------------
\item \textbf{Coeficiente de determinación}

\[
R^{2}
= \frac{SSE}{SST}
= 1-\frac{SSR}{SST}
= 1-\frac{16}{106}
\approx 0.849.
\]

Por lo tanto,
\[
\boxed{\hat\beta_0 = 4.5,\;\; \hat\beta_1 = 1.5,\;\;
SST = 106,\;\; SSR = 90,\;\; SSE = 16,\;\; R^{2}\approx 0.85.}
\]
\end{enumerate}

    \item En un curso de 7 alumnos, se tienen los siguientes datos:
    \begin{center}
\begin{tabular}{|c|c|c|c|c|}
\hline
$i$ & Nota $\, (Y_i)$ & Horas estudio $\, (X_{1i})$ & IQ $\, (X_{2i})$ & Motivación $\, (X_{3i})$\\
\hline
1 & 3 & 1 & 95 & 50\\
2 & 4 & 2 & 99 & 40\\
3 & 5 & 3 &100 & 50\\
4 & 5 & 4 &100 & 60\\
5 & 6 & 4 &105 & 70\\
6 & 7 & 5 & 99 & 70\\
7 & 7 & 6 & 95 & 70\\
\hline
\end{tabular}
\end{center}

se desea estimar la nota mediante las horas de estudio, el IQ y la motivación de la forma:
$$Nota = \beta_0 + \beta_1 X_1 + \beta_2 X_2 + \beta_3 X_3$$
\begin{enumerate}
    \item Encuentro los estimados \textbf{$\beta$}, sabiendo que la matriz $(X^tX)^{-1}$ esta dada por:
\[
(X'X)^{-1}=
\begin{pmatrix}
140.462 & -0.1628 & -1.4317 & \;\;\,0.03412\\
-0.1628 & \;\;0.201383 & \;\;0.008897 & -0.02454\\
-1.4317 & \;\;0.008897 & \;\;0.01528 & -0.001924\\
\;\;\,0.03412 & -0.02454 & -0.001924 & \;\;0.004166
\end{pmatrix}.
\]

\item Calcule $R^2$ ajustado del modelo.
\end{enumerate}

\textbf{Solución}

\paragraph{(a) Matriz de diseño y producto \(X'y\).}
\[
X=
\begin{pmatrix}
1 & 1 & 95 & 50\\
1 & 2 & 99 & 40\\
1 & 3 &100 & 50\\
1 & 4 &100 & 60\\
1 & 4 &105 & 70\\
1 & 5 & 99 & 70\\
1 & 6 & 95 & 70
\end{pmatrix},
\qquad
y=
\begin{pmatrix}3\\4\\5\\5\\6\\7\\7\end{pmatrix}.
\]

\[
X'y=
\begin{pmatrix}
\sum Y_i\\
\sum X_{1i}Y_i\\
\sum X_{2i}Y_i\\
\sum X_{3i}Y_i
\end{pmatrix}
=
\begin{pmatrix}
37\\ 147\\ 3669\\ 2260
\end{pmatrix}.
\]

\paragraph{(b) Estimador MCO.}
Con la inversa provista
\[
(X'X)^{-1}=
\begin{pmatrix}
140.462 & -0.1628 & -1.4317 & 0.03412\\
-0.1628 & 0.201383 & 0.008897 & -0.02454\\
-1.4317 & 0.008897 & 0.01528 & -0.001924\\
0.03412 & -0.02454 & -0.001924 & 0.004166
\end{pmatrix},
\]
el vector de coeficientes es
\[
\hat{\boldsymbol\beta}
=(X'X)^{-1}X'y=
\begin{pmatrix}
-2.5544\\
\;0.7668\\
\;0.04516\\
\;0.01077
\end{pmatrix}.
\]

\[
\boxed{
\hat\beta_0=-2.5544,\;
\hat\beta_1=0.7668,\;
\hat\beta_2=0.04516,\;
\hat\beta_3=0.01077}
\]

\paragraph{(c) Ecuación estimada.}
\[
\widehat{\text{Nota}}
=-2.5544
+0.7668\,(\text{horas})
+0.04516\,(\text{IQ})
+0.01077\,(\text{motivación}).
\]

--------------------------------------------------------------------
\subsection*{2.\; Cálculo de \(R^{2}\) y \(R^{2}\) ajustado}

Media de las notas  
\(\displaystyle \bar Y=\tfrac{37}{7}=5.2857\).


\begin{center}
\begin{tabular}{c|c|c|c|c|c|c}
\hline
i & Y_i & \widehat Y_i & u_i = Y_i-\widehat Y_i & Y_i-\bar Y & (Y_i-\bar Y)^2 & u_i^{\,2}\\
\hline
1 & 3 & 3.041 & -0.041 & -2.286 & 5.224 & 0.002\\
2 & 4 & 3.881 &  \;0.119 & -1.286 & 1.653 & 0.014\\
3 & 5 & 4.801 &  \;0.199 & -0.286 & 0.082 & 0.040\\
4 & 5 & 5.675 & -0.675 & -0.286 & 0.082 & 0.456\\
5 & 6 & 6.008 & -0.008 &  \;0.714 & 0.510 & 0.000\\
6 & 7 & 6.504 &  \;0.496 &  \;1.714 & 2.939 & 0.246\\
7 & 7 & 7.091 & -0.091 &  \;1.714 & 2.939 & 0.008\\
\hline
\multicolumn{5}{r|}{\textstyle\sum} & 13.429 & 0.765\\
\hline
\end{tabular}
\end{center}


\[
SST=\sum (Y_i-\bar Y)^2 = 13.4286,
\qquad
SSE=SST-SSR=12.6633.
\]

\[
R^{2}
=\frac{SSE}{SST}
=1-\frac{SSR}{SST}
=1-\frac{0.7653}{13.4286}
=0.9430.
\]

Con \(n=7\) observaciones y \(k=3\) regresores:

\[
\boxed{
\bar R^{2}=1-\frac{SSR/(n-k-1)}{SST/(n-1)}
=1-\frac{0.7653/3}{13.4286/6}
=0.8860
}.
\]

\bigskip
\noindent\textbf{Resultados finales}

\[
\boxed{
\begin{aligned}
\hat\beta_0 &= -2.5544,&
\hat\beta_1 &= 0.7668,&
\hat\beta_2 &= 0.04516,&
\hat\beta_3 &= 0.01077;\\
R^{2} &= 0.9430,&
\bar R^{2} &= 0.8860.&
\end{aligned}}
\]

\\


\item Sea un modelo de regresión lineal simple en el cual, en la práctica, se observan $
Y^* = Y + u, \quad \text{donde } u \sim \mathcal{N}(0,\,1),
$ y $
X^* = X + c$, \quad \text{donde } $c \neq 0$ \text{ es una constante}. Si el modelo verdadero (sin errores de medición) es $
Y = \beta_0 + \beta_1 \, X + \varepsilon$, muestre que, al estimar el modelo usando \((Y^*,\,X^*)\) en lugar de \((Y,\,X)\), se obtiene $E[\hat{\beta}_1^*] \;=\; \hat{\beta}_1,
\quad
E[\hat{\beta}_0^*] \;=\; \hat{\beta}_0 \;-\;\hat{\beta}_1 \, c.
$
\\
\\
\textbf{solución}
en términos esperados ,el estimador $\hat{\beta^*_1}$ esta dado por:
$$E(\hat{\beta^*_1})  =  \frac{Cov(X^*,Y^*)}{Var(X^*)} = \frac{Cov(X+c,Y+u)}{Var(X+c)}=\frac{Cov(X+c,\beta_0 + \beta_1X+u)}{Var(X+c)}$$
$$\frac{Cov(X,\beta_1X)+Cov(X,u)}{Var(X)}$$
Por supuestos de regresión lineal $Cov(X,u)=0$, por lo tanto
$$\frac{Cov(X,\beta_1X)}{Var(X)}=\frac{\beta_1Cov(X,X)}{Var(X)}=\beta_1\frac{Var(X)}{Var(X)}=\beta_1$$
y $E(\hat{\beta^*_0})$, esta dado por:
$$E(\hat{\beta^*_0})=E(Y^*)-\hat{\beta^*_1}E(X^*)=E(Y+u)-\beta_1E(X+c)$$
$$E(\hat{\beta^*_0}) = E(Y)+E(u)-\beta_1E(X)-\beta_1c$$
Se sabe que $E(u)=0$, entonces:
$$E(\hat{\beta^*_0}) = E(Y)-\beta_1E(X)-\beta_1c=\beta_0-\beta_1c$$

\textbf{solución alternativa para $\beta_1$}

$$\hat{\beta_1^*} = \frac{\sum (X_i^*-\Bar{X^*})(Y_i^*-\Bar{Y^*})}{\sum (X_i^*-\Bar{X^*})^2}$$
$$\hat{\beta_1^*} = \frac{\sum (X_i+c-(\Bar{X}+c))(Y_i+u-(\Bar{Y}+\Bar{u}))}{\sum (X_i+c-(\Bar{X}+c))^2}$$
$$\hat{\beta_1^*} = \frac{\sum (X_i-\Bar{X})(Y_i+u-(\Bar{Y}+\Bar{u}))}{\sum (X_i-\Bar{X})^2}$$
$$\hat{\beta_1^*} = \frac{\sum (X_i-\Bar{X})(Y_i-\Bar{Y})}{\sum (X_i-\Bar{X})^2} + \frac{\sum (X_i-\Bar{X})(u-\Bar{u})}{\sum (X_i-\Bar{X})^2}$$
$$\hat{\beta_1^*} = \hat{\beta_1} + \frac{\sum (X_i-\Bar{X})(u-\Bar{u})}{\sum (X_i-\Bar{X})^2}$$

el termino $\frac{\sum (X_i-\Bar{X})(u-\Bar{u})}{\sum (X_i-\Bar{X})^2}$ es igual a 0 en terminos esperados, por lo que:
 $$E(\hat{\beta_1^*}) = E(\hat{\beta_1})$$




\end{enumerate}

\section{Formulario}


\begin{itemize}

\item  Estimadores MCO para RLS.
{
$$\hat{\beta}_0 = \bar{y} - \hat{\beta}_1 \bar{x} $$
$$\hat{\beta}_1 = \frac{\sum_{i=1}^{n}(x_i - \bar{x})(y_i-\bar{y})}{\sum_{i=1}^{n}(x_i - \bar{x})^2}=\frac{Cov(x,y)}{Var(x)}$$}


\item  Estimadores MCO para RLM.
{
$${\hat{\beta}} = (X^TX)^{-1} X^TY $$}


    \item \[ R^2  = 1 - \frac{SSR}{SST} \] 
\item  $$SST = \sum_{i=1}^{n} (y_i - \bar{y})^2 \quad \quad SSR = \sum_{i=1}^{n}  \hat{u}_i^2 $$
    \item \[ \bar{R}^2 =  1 - \frac{(n-1)SSR}{(n-k-1)SST} \]

    \item  $$\hat{\sigma}^2 = \frac{\sum_{i=1}^{n} \hat{u}_i^2 }{n-k-1}= \frac{SSR}{n - k - 1}$$
    \item \[
\text{Var}(\hat{\beta}|X) = \hat{\sigma}^2(X^{\top}X)^{-1}
\]
    \end{itemize}



\end{document}